% !TEX program = xelatex
% 2027 新领军学科能力二试 · 综合 · 计算机（预测训练卷）
\documentclass[cjk=noto, seal=first]{mojing-exam}

\examtitle{2027 新领军学科能力测试 · 预测训练卷}
\secret{训练用卷 ★ 非官方试题}

\begin{document}

\begin{paper}[time={建议 1.5 小时（训练设置）}, full=100]{二试 · 综合 · 计算机}

\begin{instructions}
  \item 官方只公布二试综合六选二的科目名称，没有公布计算机的考纲、题型与分值，也未说明是否上机。本卷为笔试，
        考查算法思想、复杂度分析与信息表示，伪代码可用任何清楚的写法，全部安排为训练设置。
  \item 本卷共 4 题，每题 25 分，满分 100 分。
\end{instructions}

\begin{gradingnotes}
  \item 本评分标准为训练用建议标准。伪代码不要求特定语言，逻辑正确、边界清楚即可；分值为累计得分。
\end{gradingnotes}

\scoretable{1,2,3,4}

\begin{questions}

% =============================================================================
\question[25]
（比较的次数）只允许通过“比较两个元素的大小”获取信息。
\begin{enumerate}
  \item 在 1000 个已排好序的不同整数中查找某个数（可能不存在），二分查找在最坏情况下需要多少次比较？（每次比较可得到“小于、等于、大于”三种结果）
  \item 证明：任何基于比较的排序算法，对 $n$ 个不同元素排序，在最坏情况下至少需要 $\lceil\log_2 n!\rceil$ 次比较（这里每次比较只有“小于”“大于”两种结果；
        提示：把算法的全部可能执行过程画成一棵二叉判定树）；
  \item 分别计算 $n=5$ 和 $n=10$ 时这个下界的值；并证明 $\log_2n!\ge\frac n2\log_2\frac n2$，从而说明比较排序的最坏时间复杂度不可能低于 $n\log n$ 的量级。
\end{enumerate}
\answerspace{2cm}
\begin{solution}
(1)~每次比较后剩余候选区间至少缩小一半，$k$ 次比较最多能区分 $2^k-1$ 个位置；$2^{10}-1=1023\ge1000$，而 $2^9-1<1000$，
最坏需 10 次。\score{6}

(2)~把算法看成二叉决策树：每个内部节点是一次比较，每片叶子对应一个输出顺序。$n!$ 种输入顺序须到达不同的叶子，
所以叶子数 $\ge n!$；高为 $h$ 的二叉树至多 $2^h$ 片叶子，故 $h\ge\log_2n!$，最坏比较次数 $\ge\lceil\log_2n!\rceil$。\score{16}

(3)~$\lceil\log_2120\rceil=7$；$10!=3\,628\,800$，$2^{21}<10!<2^{22}$，下界为 22。\score{20}
由 $n!\ge(n/2)^{n/2}$，$\log_2n!\ge\frac n2\log_2\frac n2$，即下界是 $n\log n$ 量级。\score{25}
\end{solution}

% =============================================================================
\question[25]
（动态规划）用面值为 1、3、4 的硬币（每种数量不限）凑出金额 $S$。
\begin{enumerate}
  \item “每次取不超过剩余金额的最大面值”的贪心策略能否总是得到最少的硬币数？举例说明；
  \item 记 $f(S)$ 为凑出 $S$ 所需的最少硬币数，写出 $f$ 的递推关系，计算 $f(0),f(1),\dots,f(10)$，并给出求 $f(S)$ 的伪代码与时间复杂度；
  \item 不计顺序，凑出 10 共有多少种不同的方法？说明你的计算方法为什么不会把同一种组合重复计数。
\end{enumerate}
\answerspace{2cm}
\begin{solution}
(1)~不能。$S=6$ 时贪心得 $4+1+1$（3 枚），最优为 $3+3$（2 枚）。\score{5}

(2)~$f(0)=0$，$f(S)=1+\min\{f(S-c):c\in\{1,3,4\},\,c\le S\}$。依次得
$f(0..10)=0,1,2,1,1,2,2,2,2,3,3$。\score{13}
伪代码：令 $f[0]=0$；对 $s=1,\dots,S$，$f[s]=1+\min_{c\le s}f[s-c]$。共 $S$ 步，每步考察 3 种面值，时间 $O(S)$（面值有 $k$ 种时为 $O(kS)$），空间 $O(S)$。\score{17}

(3)~按面值逐种加入计数：$w[0]=1$，依次对 $c=1,3,4$，令 $s$ 从 $c$ 到 $10$ 递增执行 $w[s]\mathrel{+}=w[s-c]$，得 $w[10]=8$。
这相当于规定“先决定用几枚 1，再决定几枚 3，最后几枚 4”，每种组合只按一种顺序生成，不会重复。8 种为：
$\{4,4,1,1\},\{4,3,3\},\{4,3,1,1,1\},\{4,1^6\},\{3,3,3,1\},\{3,3,1^4\},\{3,1^7\},\{1^{10}\}$。\score{25}
\end{solution}

% =============================================================================
\question[25]
（路径计数与最短路）
\begin{enumerate}
  \item 一张由 $6\times6$ 个小正方形组成的方格网，格点坐标从 $(0,0)$ 到 $(6,6)$。从 $(0,0)$ 走到 $(6,6)$，每步只能向右或向上走一格。共有多少条路径？若格点 $(3,3)$ 不能经过，又有多少条？
  \item 写出用“广度优先搜索”求无权图中从起点到各点最短路径长度的方法，并说明它为什么正确；
  \item Dijkstra 算法求带权图从起点出发的最短路：起点的距离估计为 0，其余为 $+\infty$；每次从尚未确定的顶点中取距离估计最小者，把它的距离永久确定，
        并只用它的出边更新尚未确定顶点的距离估计。该算法要求边权非负。构造一个含负权边的三顶点有向图，说明按上述步骤会给出错误结果。
\end{enumerate}
\answerspace{2cm}
\begin{solution}
(1)~共走 12 步，选 6 步向右：$\binom{12}6=924$ 条。经过 $(3,3)$ 的有 $\binom63^2=400$ 条，故避开的有 $524$ 条。\score{8}

(2)~用队列：起点距离为 0 入队；每次取出队首 $u$，对其未访问的邻点 $v$ 令 $d(v)=d(u)+1$ 并入队。
正确性：出队的顶点按距离不减的顺序排列（队列中距离至多相差 1），某顶点第一次被访问时，由距离为 $d$ 的顶点发现，它的真实距离不可能小于 $d+1$
（否则会被更早发现），故 $d(v)$ 即最短距离。\score{17}

(3)~例：顶点 $s,a,b$，边 $s\to a$ 权 2，$s\to b$ 权 1，$a\to b$ 权 $-2$。按题述步骤先确定 $b$ 的距离为 1（当时最小），之后不再更新；
实际最短路 $s\to a\to b$ 长度为 0。负权边使“当前最近的点距离已确定”这一前提失效。\score{25}
\end{solution}

% =============================================================================
\question[25]
（纠错码）把 4 位信息编码成 7 位码字 $c=(c_1,\dots,c_7)\in\{0,1\}^7$，要求满足 $Hc\T=0$（在模 2 意义下），其中 $H$ 是 $3\times7$ 矩阵，
它的第 $j$ 列是 $j$ 的三位二进制表示（自上而下为高位到低位）：
\[
  H=\begin{pmatrix}0&0&0&1&1&1&1\\0&1&1&0&0&1&1\\1&0&1&0&1&0&1\end{pmatrix}.
\]
\begin{enumerate}
  \item 这样的码字共有多少个？证明任意两个不同码字至少有 3 位不同；
  \item 若传输中至多有 1 位出错，说明如何由接收到的 $r$ 计算 $s=Hr\T$ 来找出出错位置并纠正；
  \item 收到 $r=(1,0,0,1,0,0,0)$，假定至多 1 位出错，求原码字；
  \item 若有 2 位出错，上述方法会发生什么？
\end{enumerate}
\answerspace{2cm}
\begin{solution}
(1)~$H$ 的秩为 3（含单位列 $j=1,2,4$），解空间维数 $7-3=4$，共 $2^4=16$ 个码字。\score{4}
码是线性的，两码字之差仍是码字；若非零码字只有 1 位为 1，则 $H$ 的某一列为零，不可能；若只有 2 位为 1，则 $H$ 的两列相等，也不可能。
故非零码字至少有 3 个 1，任两码字至少 3 位不同。\score{12}

(2)~若 $r=c+e_j$（第 $j$ 位出错），则 $s=He_j\T$ 恰为第 $j$ 列，即 $j$ 的二进制；$s=0$ 表示无错。翻转第 $j$ 位即可纠正。\score{18}

(3)~$s=Hr\T$：第 1、4 位为 1，对应列 $001$ 与 $100$，和为 $101$，即 5。翻转第 5 位，原码字为 $(1,0,0,1,1,0,0)$（可验证 $Hc\T=0$）。\score{23}

(4)~若第 $i,j$ 位出错，$s$ 等于第 $i$、$j$ 列之和，即 $i$ 与 $j$ 按二进制逐位异或所得的非零数 $k\ne i,j$，算法会翻转第 $k$ 位，得到另一个码字（与原码字有 3 位不同），译码错误。最小距离为 3 的码只能纠正 1 位错误（或仅检测至多 2 位错误）。\score{25}
\end{solution}

\end{questions}
\end{paper}
\end{document}
